\documentclass[tikz,border=8pt]{standalone}
\usepackage{amsmath}
\usepackage{CJKutf8}
\usepackage{tikz}
\usetikzlibrary{arrows.meta,calc,positioning}
\begin{document}
\begin{CJK*}{UTF8}{gbsn}
\begin{tikzpicture}[
  >=Stealth,
  axis/.style={line width=1.1pt, ->},
  markline/.style={line width=0.9pt},
  callout/.style={draw=black, rounded corners=2pt, inner sep=5pt, font=\small, align=left}
]
  \node[font=\large\bfseries] at (7.1,6.35) {为什么横轴常用对数，幅值常写成 dB};

  \node[font=\small\bfseries] at (3.0,5.6) {频率轴：看差值 vs 看倍数};
  \draw[axis] (0.6,4.6) -- (5.7,4.6);
  \node[font=\scriptsize, anchor=west] at (0.6,4.95) {线性频率轴};
  \foreach \x/\lab in {0.9/{0.1},1.1/{1},1.9/{10},5.2/{100}} {
    \draw (\x,4.5) -- (\x,4.7);
    \node[font=\scriptsize] at (\x,4.28) {$\lab$};
  }
  \node[font=\scriptsize, anchor=west] at (5.9,4.6) {低频挤在一起};

  \draw[axis] (0.6,3.15) -- (5.7,3.15);
  \node[font=\scriptsize, anchor=west] at (0.6,3.5) {对数频率轴};
  \foreach \x/\lab in {0.9/{0.1},2.3/{1},3.7/{10},5.1/{100}} {
    \draw (\x,3.05) -- (\x,3.25);
    \node[font=\scriptsize] at (\x,2.83) {$\lab$};
  }
  \node[font=\scriptsize, anchor=west] at (5.9,3.15) {每一格都代表“放大 10 倍”};

  \draw[->, line width=0.9pt] (3.0,4.25) -- (3.0,3.5);
  \node[font=\scriptsize] at (3.55,3.88) {从看差值变成看倍数};

  \node[font=\small\bfseries] at (10.2,5.6) {幅值：乘法关系改写成 dB};
  \draw[markline] (8.2,2.6) -- (8.2,5.0);
  \draw[markline] (7.9,4.5) -- (8.5,4.5);
  \draw[markline] (7.9,3.8) -- (8.5,3.8);
  \draw[markline] (7.9,3.1) -- (8.5,3.1);
  \node[font=\scriptsize, anchor=east] at (7.8,4.5) {$+20\ \mathrm{dB}$};
  \node[font=\scriptsize, anchor=east] at (7.8,3.8) {$0\ \mathrm{dB}$};
  \node[font=\scriptsize, anchor=east] at (7.8,3.1) {$-20\ \mathrm{dB}$};

  \node[callout, anchor=west] at (8.8,4.5) {幅值 $10$ 倍\\写成 $+20\,\mathrm{dB}$};
  \node[callout, anchor=west] at (8.8,3.8) {幅值 $1$ 倍\\写成 $0\,\mathrm{dB}$};
  \node[callout, anchor=west] at (8.8,3.1) {幅值 $0.1$ 倍\\写成 $-20\,\mathrm{dB}$};

  \node[callout, anchor=north west] at (8.4,2.35) {读图时先抓两件事：\\1. 横轴看频率放大了多少倍\\2. 纵轴看幅值相差了多少倍};
\end{tikzpicture}
\end{CJK*}
\end{document}
